\documentclass[conference]{IEEEtran}
\usepackage[utf8]{inputenc}
\usepackage{graphicx}
\usepackage{amsmath}
\usepackage{cite}
\usepackage{hyperref}
\usepackage{listings}
\usepackage{xcolor}
\usepackage{booktabs}
\usepackage{tabularx}

% Code listing style
\lstset{
    basicstyle=\small\ttfamily,
    breaklines=true,
    frame=single,
    numbers=left,
    numberstyle=\tiny,
    keywordstyle=\color{blue},
    commentstyle=\color{gray},
    stringstyle=\color{red},
    language=C
}

\begin{document}

\title{Alarm and Reminder System: A Miniature Task Scheduler Demonstrating OS-Level Time-Based Event Management and Process Triggers}

\author{
    \IEEEauthorblockN{Aman Kumar}
    \IEEEauthorblockA{
        \\
        B.Tech Computer Science and AI \\
    }
}

\maketitle

\begin{abstract}
This paper presents the design and implementation of an Alarm and Reminder System that functions as a miniature operating system task scheduler. The system demonstrates fundamental OS concepts including process creation via \texttt{fork()}, background task execution through kernel-managed \texttt{sleep()}, signal-based process termination using \texttt{kill()}, timezone handling via process environment manipulation, and zombie prevention through \texttt{SA\_NOCLDWAIT}. Built using a three-tier architecture comprising a C alarm engine, a Node.js bridge server, and a web-based frontend, the system mirrors the behavior of real-world task schedulers such as cron and systemd timers. Each alarm is scheduled as an independent child process---created via \texttt{fork()}, suspended via the kernel's process scheduler during \texttt{sleep()}, and awakened by timer interrupts to execute the alarm action. The system implements duplicate alarm prevention through PID file-based process tracking, safe delay computation using \texttt{mktime()} and \texttt{difftime()}, and cross-platform sound playback via \texttt{execlp()}. Testing confirms accurate future time calculation, safe process management, and proper resource cleanup. This project bridges the gap between theoretical OS concepts and practical implementation, providing a tangible demonstration of how operating systems manage timed events and process lifecycle.
\end{abstract}

\begin{IEEEkeywords}
Operating Systems, Process Scheduling, fork(), Task Management, Signal Handling
\end{IEEEkeywords}

\section{Introduction}

Operating systems manage the execution of processes, memory allocation, and hardware interaction, forming the foundation of all computing. Among the most critical OS responsibilities is \textbf{task scheduling}---determining when and how processes execute, managing their lifecycle from creation to termination, and handling time-based triggers that automate system operations \cite{silberschatz2018}.

Real-world systems rely heavily on timed events: cron jobs run periodic maintenance scripts, \texttt{at(1)} schedules one-time tasks, systemd timers manage service lifecycles, and the kernel itself uses timer interrupts to preempt processes and maintain system responsiveness \cite{tanenbaum2015}. Despite their ubiquity, these mechanisms are often taught abstractly, leaving students without practical experience building systems that use them.

This project addresses that gap by implementing an \textbf{Alarm and Reminder System} that functions as a miniature task scheduler. Rather than merely calling library functions, the system demonstrates the complete lifecycle of a scheduled task:

\begin{enumerate}
    \item \textbf{Parsing and validation} of human-readable time specifications
    \item \textbf{Conversion} to UNIX timestamps using \texttt{mktime()}
    \item \textbf{Delay computation} using \texttt{difftime()} for safe scheduling
    \item \textbf{Process creation} via \texttt{fork()} to isolate each alarm
    \item \textbf{Kernel-managed suspension} via \texttt{sleep()} using timer interrupts
    \item \textbf{Background execution} independent of the parent process
    \item \textbf{Signal-based cancellation} using \texttt{kill(pid, SIGTERM)}
    \item \textbf{Resource cleanup} to prevent zombie processes
\end{enumerate}

The system also demonstrates multi-timezone time handling, cross-platform sound playback, and inter-process communication through file-based PID tracking---concepts that are foundational to OS design but rarely combined in a single educational project.

\section{Literature Review}

\subsection{Process Scheduling in Operating Systems}

Silberschatz et al. \cite{silberschatz2018} define process scheduling as the activity of the OS kernel that selects which process should execute next. Modern schedulers like the Completely Fair Scheduler (CFS) on Linux use red-black trees to maintain $O(\log n)$ scheduling decisions, while macOS employs a multilevel feedback queue \cite{love2010}. Both rely on timer interrupts to preempt running processes and enforce time-sharing.

The \texttt{fork()} system call, introduced in Unix V1 (1971), remains the primary mechanism for process creation across POSIX-compliant systems \cite{stevens2013}. When \texttt{fork()} is called, the kernel creates a new process with an independent PID, copies the parent's page table entries (using copy-on-write for efficiency), and returns control to both parent and child with different return values.

\subsection{Time Management and the UNIX Epoch}

POSIX time management centers on the UNIX epoch (January 1, 1970, 00:00:00 UTC) \cite{ieee2017}. The \texttt{time()} system call reads the kernel's system clock, returning seconds since the epoch. The \texttt{mktime()} function converts broken-down local time (\texttt{struct tm}) to a \texttt{time\_t} timestamp, accounting for timezone and daylight saving time.

\subsection{Signal Handling and Process Control}

Stevens and Rago \cite{stevens2013} describe signals as software interrupts delivered to processes by the kernel. \texttt{SIGTERM} (signal 15) requests graceful termination, while \texttt{SIGCHLD} notifies parents of child state changes. The \texttt{SA\_NOCLDWAIT} flag instructs the kernel to automatically reap terminated children, preventing resource leaks from zombie processes \cite{kerrisk2010}.

\subsection{Educational Approaches to OS Concepts}

Tanenbaum and Bos \cite{tanenbaum2015} advocate for building small operating systems (MINIX) to teach OS concepts. Similarly, xv6 (MIT) \cite{cox2019} provides a simplified Unix implementation. This project takes a complementary approach---building a user-space application that extensively uses OS system calls.

\subsection{Task Scheduling Systems}

The cron daemon remains the standard for periodic task scheduling \cite{ieee2017}. Modern alternatives include systemd timers and launchd. All follow the same pattern: parse a time specification, compute the delay, create a process, and execute the task. Our system replicates this pattern at a smaller scale.

\section{Methodology}

\subsection{System Architecture}

The system employs a three-tier architecture:

\textbf{Layer 1---C Alarm Engine:} The core scheduler directly interfaces with the OS kernel through system calls including \texttt{fork()}, \texttt{mktime()}, \texttt{difftime()}, \texttt{sleep()}, \texttt{kill()}, \texttt{setenv()}, and \texttt{execlp()}.

\textbf{Layer 2---Node.js Bridge Server:} An Express.js HTTP server that uses \texttt{child\_process.execFile()} to invoke the C engine, demonstrating higher-level process management.

\textbf{Layer 3---Web Frontend:} A browser-based interface providing alarms, reminders, and world clock functionality via HTTP REST APIs.

\subsection{Alarm Scheduling Algorithm}

The scheduling process follows this sequence:

\begin{lstlisting}[caption={Alarm Scheduling Algorithm},language=C]
schedule_alarm(datetime, message, sound):
  1. Parse datetime -> struct tm
  2. target = mktime(struct tm)
  3. now = time(NULL)
  4. delay = difftime(target, now)
  5. If delay <= 0: REJECT (past time)
  6. If is_duplicate(): REJECT
  7. Set SA_NOCLDWAIT on SIGCHLD
  8. pid = fork()
     Parent: register_pid(), return JSON
     Child: sleep(delay), play_sound(),
            unregister_pid(), _exit(0)
\end{lstlisting}

\subsection{Duplicate Prevention Mechanism}

The system maintains a PID tracking file at \texttt{/tmp/alarm\_engine\_pids.dat}. Before scheduling, the engine sends a null signal \texttt{kill(pid, 0)} to verify each existing process is alive. If an alive process matches the same datetime and message, the request is rejected. Dead entries are automatically cleaned up.

\subsection{Timezone Handling}

The world clock queries 8 IANA timezones by:
\begin{enumerate}
    \item Calling \texttt{setenv("TZ", timezone, 1)} to modify the process environment
    \item Calling \texttt{tzset()} to re-read the TZ variable
    \item Calling \texttt{localtime(\&now)} to interpret in the new timezone
    \item Extracting formatted time, date, and UTC offset
\end{enumerate}

\subsection{Cross-Platform Sound Playback}

Sound playback uses OS detection:
\begin{itemize}
    \item \textbf{macOS:} \texttt{access()} for file check, then \texttt{fork()} + \texttt{execlp("afplay")}
    \item \textbf{Linux:} \texttt{access("/usr/bin/paplay")} then \texttt{execlp("paplay")}
    \item \textbf{Fallback:} Terminal BEL character \texttt{$\backslash$a}
\end{itemize}

\section{Results and Discussion}

\subsection{Functional Verification}

All system features were tested and verified as shown in Table~\ref{tab:results}.

\begin{table}[h]
\centering
\caption{Feature Testing Results}
\label{tab:results}
\begin{tabular}{lll}
\toprule
\textbf{Feature} & \textbf{Test Method} & \textbf{Result} \\
\midrule
Alarm scheduling & CLI + API & Pass \\
World clock & 8 timezones & Pass \\
Sound playback & test-sound cmd & Pass \\
Task cancellation & Cancel by PID & Pass \\
Duplicate prevention & Double schedule & Pass \\
Past time rejection & Past datetime & Pass \\
Zombie prevention & ps inspection & Pass \\
Active alarm listing & list command & Pass \\
\bottomrule
\end{tabular}
\end{table}

\subsection{Process Lifecycle Observation}

Using \texttt{ps -o pid,ppid,stat,command}, we observed:
\begin{enumerate}
    \item After \texttt{fork()}, the child appears with state \textbf{S} (sleeping)
    \item During the delay, the child consumes 0\% CPU
    \item At alarm time, the child transitions to \textbf{R} (running)
    \item With \texttt{SA\_NOCLDWAIT}, no zombie (state Z) appears
\end{enumerate}

\subsection{Performance Characteristics}

\begin{itemize}
    \item \textbf{Process creation:} \texttt{fork()} completes in $< 1$ ms (copy-on-write)
    \item \textbf{Memory per alarm:} $\approx 1$ MB RSS per sleeping child
    \item \textbf{Timer accuracy:} Within 1 second of target (limited by \texttt{sleep()} granularity)
\end{itemize}

\subsection{Comparison with Real Schedulers}

Table~\ref{tab:comparison} compares our system with production schedulers.

\begin{table}[h]
\centering
\caption{Comparison with Production Schedulers}
\label{tab:comparison}
\begin{tabular}{llll}
\toprule
\textbf{Aspect} & \textbf{Ours} & \textbf{cron} & \textbf{systemd} \\
\midrule
Mechanism & fork+sleep & fork+exec & service \\
Persistence & PID file & crontab & .timer \\
Granularity & 1 sec & 1 min & 1 $\mu$s \\
Cancel & SIGTERM & crontab -r & systemctl \\
\bottomrule
\end{tabular}
\end{table}

\section{Conclusion and Future Work}

\subsection{Conclusion}

This project successfully demonstrates that alarms and reminders are fundamentally OS-level scheduling problems. By implementing a three-tier system using \texttt{fork()}, \texttt{sleep()}, \texttt{mktime()}, \texttt{difftime()}, \texttt{kill()}, and signal handling, we show how operating systems manage timed events and process triggers.

The key contributions are: (1) a working miniature task scheduler mirroring cron's behavior; (2) visible OS concepts with documented system calls; (3) complete process lifecycle management; (4) practical features including world clock, sound playback, and cancellation; and (5) professional tooling with Makefile, debugger support, and process inspection.

\subsection{Future Work}

Future enhancements include: replacing \texttt{sleep()} with POSIX \texttt{timer\_create()} for kernel-level timers; using \texttt{mmap()} with \texttt{MAP\_SHARED} for shared memory IPC; implementing \texttt{pthread\_create()} to compare threading vs. forking; adding \texttt{flock()} for file locking; and exploring Linux namespaces for container isolation.

\begin{thebibliography}{10}

\bibitem{silberschatz2018}
A.~Silberschatz, P.~B.~Galvin, and G.~Gagne, \emph{Operating System Concepts}, 10th~ed.\hskip 1em plus 0.5em minus 0.4em Hoboken, NJ, USA: Wiley, 2018.

\bibitem{tanenbaum2015}
A.~S.~Tanenbaum and H.~Bos, \emph{Modern Operating Systems}, 4th~ed.\hskip 1em plus 0.5em minus 0.4em Upper Saddle River, NJ, USA: Pearson, 2015.

\bibitem{love2010}
R.~Love, \emph{Linux Kernel Development}, 3rd~ed.\hskip 1em plus 0.5em minus 0.4em Upper Saddle River, NJ, USA: Addison-Wesley, 2010.

\bibitem{stevens2013}
W.~R.~Stevens and S.~A.~Rago, \emph{Advanced Programming in the UNIX Environment}, 3rd~ed.\hskip 1em plus 0.5em minus 0.4em Upper Saddle River, NJ, USA: Addison-Wesley, 2013.

\bibitem{ieee2017}
IEEE, ``IEEE Std 1003.1-2017---Standard for Information Technology---Portable Operating System Interface (POSIX),'' The Open Group, 2017.

\bibitem{kerrisk2010}
M.~Kerrisk, \emph{The Linux Programming Interface}.\hskip 1em plus 0.5em minus 0.4em San Francisco, CA, USA: No Starch Press, 2010.

\bibitem{cox2019}
R.~Cox, M.~F.~Kaashoek, and R.~Morris, ``xv6, a simple Unix-like teaching operating system,'' MIT PDOS, 2019.

\bibitem{bovet2005}
D.~P.~Bovet and M.~Cesati, \emph{Understanding the Linux Kernel}, 3rd~ed.\hskip 1em plus 0.5em minus 0.4em Sebastopol, CA, USA: O'Reilly, 2005.

\bibitem{kernighan1988}
B.~W.~Kernighan and D.~M.~Ritchie, \emph{The C Programming Language}, 2nd~ed.\hskip 1em plus 0.5em minus 0.4em Englewood Cliffs, NJ, USA: Prentice Hall, 1988.

\bibitem{bach1986}
M.~J.~Bach, \emph{The Design of the UNIX Operating System}.\hskip 1em plus 0.5em minus 0.4em Englewood Cliffs, NJ, USA: Prentice Hall, 1986.

\end{thebibliography}

\end{document}
